% ============================================================
% 表4.X: 6艘代表船GA-RH实验结果
% 插入位置：论文第四章 实验与分析
% ============================================================
\begin{table}[htbp]
\centering
\caption{6艘代表外贸船GA-RH配载优化实验结果}
\label{tab:6ship_ga_rh}
\small
\begin{tabular}{lccccccccc}
\toprule
\textbf{船代码} & \textbf{箱数} & \textbf{POD数} & \textbf{Fitness} & $f_1$\textsuperscript{a} & $f_2$\textsuperscript{b} & $f_3$\textsuperscript{c} & $f_4$\textsuperscript{d} & \textbf{罚分} & \textbf{耗时(s)} \\
\midrule
CNTIG & 583 & 2 & 0.7537 & 0.9368 & 0.7786 & 0.9881 & 1.0000 & 0.0300 & 64 \\
CNCT  & 795 & 4 & 0.6999 & 0.7491 & 0.7163 & 0.9884 & 1.0000 & 0.0270 & 288 \\
MXNT  & 1,492 & 5 & 0.6201 & 0.6239 & 0.5534 & 0.9978 & 1.0000 & 0.0258 & 813 \\
OFUT  & 1,761 & 7 & 0.5900 & 0.6191 & 0.5471 & 0.9827 & 1.0000 & 0.0304 & 1,050 \\
CGAMV & 2,993 & 8 & 0.4854 & 0.6043 & 0.5994 & 0.9908 & 1.0000 & 0.0546 & 3,012 \\
APESP & 4,008 & 9 & 0.4592 & 0.5802 & 0.5438 & 0.9932 & 1.0000 & 0.0549 & 2,459 \\
\bottomrule
\end{tabular}
\vspace{2mm}
\parbox{0.9\textwidth}{\footnotesize
\textsuperscript{a}$f_1$：翻箱成本（式4.10），值越高翻箱越少；\\
\textsuperscript{b}$f_2$：同港集中效率（式4.11），值越高集中度越好；\\
\textsuperscript{c}$f_3$：重量均衡度（式4.12），值越高各列重量越均衡；\\
\textsuperscript{d}$f_4$：堆场协同度（式4.13），值越高与堆场推荐位置匹配越好。
}
\end{table}
